% Section 1: The dual space: covectors} | Part: 1 | Depends on: see docs/REWRITE_PLAN.md | Status: todo | Original lines: 398-452
%!TEX root = ../../../main.tex
\subsection{The dual space: covectors}\label{sec:dualspace}

\begin{formal}[{dual space and dual basis}]
The \emph{dual space} of $E$ is the set of all linear maps from $E$ to $\R$:
\[
E^*=\{f:E\to\R\ \text{linear}\}.
\]
It is a vector space of dimension $m$; its elements are called \emph{covectors} (or linear forms, or 1-forms). Given a basis $\{e_i\}$ of $E$ there is a unique \emph{dual basis} $\{e^{i}\}$ of $E^*$ determined by
\begin{equation}\label{eq:dualbasis}
e^{i}(e_j)=\delta^i{}_j\qquad(\text{1 if }i=j,\ 0\text{ otherwise}).
\end{equation}
Consequences:
\begin{itemize}[nosep]
\item $e^{i}(v)=v^i$: the dual basis \emph{extracts components}.
\item Every covector is $f=f_i\,e^{i}$ with $f_i=f(e_i)$, and $f(v)=f_i\,v^i$.
\item $(E^*)^*\cong E$ canonically (finite dimension): a vector $v$ is a linear machine that eats a covector $f$ and returns $f(v)$.
\end{itemize}
\end{formal}

\begin{intuition}[{covectors are ``measuring devices''}]
A vector is \emph{something that you measure}; a covector is \emph{something that measures}. It eats an arrow and returns a number, linearly. Picture a covector as a \textbf{stack of equally spaced parallel hyperplanes} (level sets of $f$): $f(v)$ is the number of layers that the arrow $v$ pierces. A short-spaced stack is a ``large'' covector.

\emph{Examples.} A ruler's coordinate function $v\mapsto v^i$; a thermometer-gradient $v\mapsto$ (change in temperature along $v$); a ``price per unit'' functional $v\mapsto$ cost of the shopping vector $v$. Covectors are \emph{not} arrows. They only become arrows once a metric is given (Section~\ref{sec:metric}).
\end{intuition}

\subsection*{Change of basis in the dual}
If $e'_i=A^j{}_ie_j$, the new dual basis is $e'^{i}=(A^{-1})^i{}_je^{j}$ (check with \eqref{eq:dualbasis}), and the components of a covector change as
\begin{equation}\label{eq:covchange}
f'_i=f(e'_i)=A^j{}_i\,f_j .
\end{equation}
In matrix language the dual basis changes with the \emph{inverse transpose}: $A^{-t}$, and the components with $A^{t}$.

\begin{intuition}[{covariant = changes like the basis}]
Covector components transform with the \emph{same} matrix $A$ as the basis vectors, so they are \textbf{covariant}. Doubling the basis vector $e_1$ doubles $f(e_1)=f_1$. Components of vectors $\leftrightarrow$ opposite of the basis; components of covectors $\leftrightarrow$ same as the basis. This is the mechanism that makes $f(v)=f_iv^i$ basis independent:
$f'_iv'^i=A^j{}_if_j\,(A^{-1})^i{}_kv^k=f_jv^j$.
\end{intuition}

\begin{example}[{a dual basis is not ``the coordinates''}]
In $\R^2$ take the basis $e_1=(1,0)$, $e_2=(1,1)$. We look for $e^{1},e^{2}$ with $e^{i}(e_j)=\delta^i_j$. Writing $e^{1}(x,y)=ax+by$ we need $a=1$ and $a+b=0$, so $e^{1}(x,y)=x-y$; similarly $e^{2}(x,y)=y$.

Notice that $e^{1}\neq$ ``the first Cartesian coordinate'': \textbf{the dual vector $e^{1}$ depends on the whole basis}, not only on $e_1$. If you only change $e_2$ then $e^{1}$ changes! Check with $v=(3,1)=2e_1+1e_2$: $e^{1}(v)=3-1=2$, $e^{2}(v)=1$. \checkmark
\end{example}

\begin{example}[{the differential of a function is a covector}]
Let $f:\R^2\to\R$, $f(x,y)=x^2y$. At $p=(1,2)$, $\dd f_p(v)=\frac{d}{d\varepsilon}f(p+\varepsilon v)|_0=\partial_xf\,v^1+\partial_yf\,v^2=4v^1+1\,v^2$. So $\dd f_p=4\,\dd x+1\,\dd y$ is a covector (the \emph{gradient} before using a metric). Its ``stack of lines'' is the family of tangent lines to the level curves of $f$.
\end{example}

\begin{pitfall}[{gradient is a covector}]
In school you drew the gradient as an arrow. That arrow secretly used the Euclidean dot product to turn the covector $\dd f$ into a vector. In curved spacetime we will always have to say which metric does this, see Section~\ref{sec:metric}.
\end{pitfall}
